% crosspolytope.tex — \input'd from experiments.tex
%
% EHZ capacity of the 4D crosspolytope (hyperoctahedron, 16 facets).
%
% Sources:
%   - Computed data: experiments/crosspolytope/crosspolytope.jsonl
%   - Binary: experiments/crosspolytope/crosspolytope.rs
%   - Symmetry group: analytical + Rust enumeration (match confirmed)
%
% Structure:
%   - Setup and motivation
%   - Symmetry reduction
%   - Result and comparison with hypercube

\subsection{Crosspolytope Capacity}
\label{sec:crosspolytope}

The 4D crosspolytope~$K^{\diamond} = \operatorname{conv}\{\pm e_1, \pm e_2,
\pm e_3, \pm e_4\}$ is the dual of the hypercube~$[-1,1]^4$.
It has~$16$ facets (one per sign vector in~$\{\pm 1\}^4$), $8$~vertices,
and is non-simple: each vertex lies on~$8$ facets.
No literature value for~$c_{\mathrm{EHZ}}(K^{\diamond})$ is known to us.

% [TODO: JÖRN - verify the symmetry group computation below. Agent computed
%  analytically and by Rust enumeration (both give 32). Check the argument.]

\subsubsection*{Symmetry reduction}

The hyperoctahedral group~$\operatorname{Aut}(K^{\diamond})$ has order~$384 =
2^4 \cdot 4!$ (signed permutations of coordinates).
For the capacity computation, only the symplectic subgroup
$G = \operatorname{Aut}(K^{\diamond}) \cap \mathrm{Sp}(4,\mathbb{R})$
acts on Reeb orbits. We computed~$|G| = 32$ both by enumeration over
all~$384$ signed permutation matrices and analytically:
a coordinate permutation~$\pi$ preserves~$\omega_0$ if and only if it maps
the symplectic planes~$\{(q_1,p_1)\}$ and~$\{(q_2,p_2)\}$ to symplectic
planes, giving~$8$ valid permutations. For each, symplecticity constrains
the signs to~$2$ free choices, for a total of~$8 \times 4 = 32$.

Using~$G$, we replace each subset~$S$ of facets by its canonical
representative, reducing the number of subsets to examine by a factor
of~${\sim}28\times$ across subset sizes.

\subsubsection*{Result}

We searched all subset sizes~$m = 2, \ldots, 12$ exhaustively
($12.2 \times 10^6$ adjacent cyclic permutations evaluated,
approximately~$6$~minutes in release mode).
Subset sizes~$m = 13, \ldots, 16$ were not computed: the permutation
counts grow factorially, and the minimum-action orbit appears at~$m = 4$.

\begin{table}[h]
\centering
\begin{tabular}{ll}
\toprule
$c_{\mathrm{EHZ}}(K^{\diamond})$ & $4.0$ \\
$\operatorname{vol}_4(K^{\diamond})$ & $\tfrac{32}{3} \approx 10.667$ \\
$\mathrm{sys}(K^{\diamond})$ & $\tfrac{3}{4} = 0.75$ \\
Minimising orbit & $m = 4$, $\beta_i = \tfrac{1}{4}$ for all~$i$ \\
\bottomrule
\end{tabular}
\caption{EHZ capacity and systolic ratio of the 4D crosspolytope.}
\label{tab:crosspolytope}
\end{table}

The crosspolytope satisfies Viterbo's conjecture with~$\mathrm{sys} =
\tfrac{3}{4}$.
Its capacity~$c_{\mathrm{EHZ}}(K^{\diamond}) = 4.0$ equals that of its dual,
the hypercube~$[-1,1]^4$.
The systolic ratios differ ($\tfrac{3}{4}$ vs.\ $\tfrac{1}{2}$) because
the volumes differ ($\tfrac{32}{3}$ vs.\ $16$).

% [TODO: JÖRN - is c_EHZ(K) = c_EHZ(K°) expected for dual polytopes?
%  We have one data point (hypercube/crosspolytope). If this is a known
%  result or has a simple explanation, cite it. If not, note as observation.]
